% =====================================================================
% Wave file: Kapranov L_infty structure on T_X[-1] via iterated Atiyah
% classes, specialised to CY_3. Three Atiyah cocycles: At(T_X), tree-level
% Yukawa Y_3, one-loop BCOV curving alpha_BCOV -- chain level.
%
% Side-by-side with Costello--Francis--Gwilliam 2026 (CFG 2026):
% E_n factorisation algebras, formality, anomaly-tree control.
%
% Standalone .tex; NOT included in main.tex. Do not read as manuscript.
% =====================================================================

\documentclass[11pt]{article}
\usepackage[localtheorems]{raeez-math-template}

\newtheorem{theorem}{Theorem}[section]
\newtheorem{proposition}[theorem]{Proposition}
\newtheorem{lemma}[theorem]{Lemma}
\newtheorem{corollary}[theorem]{Corollary}
\newtheorem{definition}[theorem]{Definition}
\newtheorem{remark}[theorem]{Remark}
\newtheorem{question}[theorem]{Question}
\newtheorem{attack}[theorem]{Attack}
\newtheorem{heal}[theorem]{Heal}
\newtheorem{surviving}[theorem]{Surviving core}

\providecommand{\At}{\mathrm{At}}
\providecommand{\End}{\mathrm{End}}
\providecommand{\CY}{\mathrm{CY}}
\providecommand{\Ethree}{E_3}
\providecommand{\Etwo}{E_2}
\providecommand{\Eone}{E_1}
\providecommand{\BCOV}{\mathrm{BCOV}}
\providecommand{\hCS}{\mathrm{hCS}}
\providecommand{\HH}{\mathrm{HH}}
\providecommand{\Coh}{\mathrm{Coh}}
\providecommand{\Perf}{\mathrm{Perf}}
\providecommand{\Om}{\Omega}
\providecommand{\Sym}{\mathrm{Sym}}
\providecommand{\id}{\mathrm{id}}
\providecommand{\cA}{\mathcal{A}}
\providecommand{\cC}{\mathcal{C}}
\providecommand{\cE}{\mathcal{E}}
\providecommand{\cF}{\mathcal{F}}
\providecommand{\cG}{\mathcal{G}}
\providecommand{\cM}{\mathcal{M}}
\providecommand{\cO}{\mathcal{O}}
\providecommand{\cT}{\mathcal{T}}
\providecommand{\cH}{\mathcal{H}}
\providecommand{\fg}{\mathfrak{g}}
\providecommand{\RR}{\mathbb{R}}
\providecommand{\CC}{\mathbb{C}}
\providecommand{\NN}{\mathbb{N}}
\providecommand{\ZZ}{\mathbb{Z}}
\providecommand{\kBKM}{\kappa_{\mathrm{BKM}}}
\providecommand{\kch}{\kappa_{\mathrm{ch}}}
\providecommand{\kcat}{\kappa_{\mathrm{cat}}}
\providecommand{\kfib}{\kappa_{\mathrm{fiber}}}
\providecommand{\Conf}{\mathrm{Conf}}
\providecommand{\MC}{\mathrm{MC}}
\providecommand{\Sp}{\mathrm{Sp}}
\providecommand{\FM}{\overline{\mathrm{FM}}}

\title{Kapranov's $L_\infty$-structure on $T_X[-1]$ via iterated Atiyah
classes, specialised to $\CY_3$: chain-level attack--heal cycles on the
three Atiyah cocycles, side-by-side with Costello--Francis--Gwilliam
2026}
\author{(Wave file; not part of the manuscript)}

\begin{document}

\maketitle

\section*{Orientation}

The $\cO_X$-linear Atiyah class $\At(T_X) \in H^1(X, \Om^1_X \otimes \End T_X)$
of the holomorphic tangent bundle is a single cohomology class; yet when
it is exhibited as the binary bracket of Kapranov's $L_\infty$-structure on
$T_X[-1]$, it sources at least two further, cohomologically independent,
classes in the $\CY_3$ Hodge diamond: the tree-level Yukawa $Y_3 \in
H^{0,3}(X) \simeq \Sym^3 H^1(X,T_X)^\vee$ and the one-loop BCOV curving
$\alpha_\BCOV = (\chi(X)/24)[\Om_X]^{0,1} \in H^{0,1}(X)$. This file
dissects the chain-level derivation of all three, re-attacks the
statement five times, and heals it five times, always with explicit
chain-level witnesses and with side-by-side comparison to the Costello--
Francis--Gwilliam 2026 $E_n$-factorisation framework. The three classes
are (i) the Kapranov generator (degree~$1$ in $H^1(X,\Om^1 \otimes \End
T_X)$), (ii) the tree-level Yukawa sourcing the cubic $\hCS$ action, and
(iii) the one-loop BV anomaly of $\hCS$. They lie in Hodge-disjoint
bidegrees and feed Stage~1 $\PhiFA_3$ at orders $\hbar^0$, tree, and
$\hbar^1$-loop respectively.

\tableofcontents

% =====================================================================
\section{Foundational material: Atiyah class, Kapranov $L_\infty$, and
the three Hodge slots}
\label{sec:foundations}

\subsection{The Atiyah class from first principles}
\label{subsec:atiyah-first-principles}

\begin{definition}[Atiyah class of a holomorphic vector bundle]
\label{def:atiyah-class}
Let $X$ be a complex manifold and $E \to X$ a holomorphic vector bundle.
The $1$-jet exact sequence
\[
  0 \longrightarrow \Om^1_X \otimes E \longrightarrow J^1(E)
  \longrightarrow E \longrightarrow 0
\]
of holomorphic sheaves is classified, in $\mathrm{Ext}^1_{\cO_X}(E,
\Om^1_X \otimes E)$, by the Atiyah class
\[
  \At(E) \in \mathrm{Ext}^1_{\cO_X}(E, \Om^1_X \otimes E)
  \simeq H^1(X, \Om^1_X \otimes \End E).
\]
Equivalently, choosing any smooth $(1,0)$-connection $\nabla^{(1,0)}$ on
$E$, the Dolbeault representative is
\begin{equation}
\label{eq:atiyah-dolbeault}
  \At(E) = [\bar\partial, \nabla^{(1,0)}]
  \in A^{0,1}(X, \Om^1_X \otimes \End E).
\end{equation}
The cohomology class is independent of the choice of $\nabla^{(1,0)}$.
\end{definition}

\begin{remark}[What the Atiyah class measures]
Cohomologically, $\At(E)$ is the obstruction to the existence of a
\emph{holomorphic} connection on $E$. It is the holomorphic analogue of
the Chern curvature $F_\nabla$ of a Hermitian connection, but without the
reality constraint: one keeps the Dolbeault $(1,1)$-component
$[\bar\partial, \nabla^{(1,0)}]$ and discards the $(2,0)$-component
$(\nabla^{(1,0)})^2$ because the latter is a tensor in the algebra of
polyvector fields and does not sit in the Atiyah slot. In the
Dolbeault-HKR model used throughout the chain-level programme, $\At(E)$
is an element of $A^{0,1}(X, \Om^1_X \otimes \End E)$ whose
$\bar\partial$-class realises the obstruction.
\end{remark}

\begin{proposition}[Local Dolbeault formula]
\label{prop:local-dolbeault}
On a coordinate patch $U \subset X$ with holomorphic coordinates
$(z^a)$ and unitary frame $\{e_i\}$ for $E|_U$, the Dolbeault Atiyah
representative reads
\begin{equation}
\label{eq:local-dolbeault}
  \At(E)|_U
  = -\bar\partial(\Gamma^i_{ja}) \, dz^a \otimes dz^j \otimes (e_i \otimes e^j)
  \in A^{0,1}(U, \Om^1 \otimes \End E),
\end{equation}
where $\Gamma^i_{ja}$ are the Christoffel symbols of the smooth
connection $\nabla^{(1,0)} e_j = \Gamma^i_{ja} \, dz^a \, e_i$. In
particular, for $E = T_X$ with a K\"ahler metric,
\begin{equation}
\label{eq:atiyah-kaehler}
  \At(T_X) = R^i_{jk\bar\ell}\, dz^k \otimes d\bar z^\ell \otimes dz^j
  \otimes \partial_i,
\end{equation}
where $R^i_{jk\bar\ell}$ is the Chern curvature of the Levi-Civita
connection.
\end{proposition}

\begin{proof}
Specialise the $A^{0,1}$ component of $[\bar\partial, \nabla^{(1,0)}]$
to the Levi-Civita case and unpack the bundle index structure.
\end{proof}

\begin{remark}[Relation to Chern classes]
Under the invariant polynomial trace, $\mathrm{tr}(\At(E)^k)/k! = c_k(E)$
in $H^k(X, \Om^k_X) \subset H^{2k}(X, \CC)$, giving a Dolbeault
representative of the Chern character. For $E = T_X$ a $\CY_n$ tangent
bundle, $c_1(T_X) = 0$ in $H^1(X, \Om^1_X)$, so the cohomology trace of
$\At(T_X)$ vanishes; yet $\At(T_X)$ itself does not vanish, because its
$\End T_X$-trace is only one of many scalar invariants.
\end{remark}

\subsection{The Kapranov $L_\infty$-structure on $T_X[-1]$}
\label{subsec:kapranov-linf}

\begin{theorem}[Kapranov 1999, Compositio 115]
\label{thm:kapranov}
Let $X$ be a complex manifold and $T_X$ the holomorphic tangent bundle.
The shifted Dolbeault resolution
\[
  \mathfrak g_X := A^{0, \bullet}(X, T_X[-1])
\]
carries a canonical $L_\infty$-structure $\{\ell_n\}_{n \geq 1}$, natural
in $X$, with the following properties:
\begin{enumerate}
\item $\ell_1 = \bar\partial$ is the Dolbeault differential.
\item $\ell_2(v,w) = \At(T_X) \cdot (v \wedge w)$, antisymmetrised, is
the binary bracket whose cohomology class is the Atiyah class $\At(T_X)$
contracted through $\End T_X$.
\item $\ell_n$ for $n \geq 3$ is given by an iterated-Atiyah formula
\[
  \ell_n(v_1, \ldots, v_n) = \sum_T w_T \cdot \mathrm{Iterate}_T(\At(T_X))(v_1,
    \ldots, v_n),
\]
where the sum runs over binary rooted trees $T$ with $n$ leaves, $w_T$
are rational weights, and $\mathrm{Iterate}_T(\At(T_X))$ is the iterated
composition of $\At(T_X)$ along $T$.
\item The $L_\infty$-relations $\sum_{p+q=n+1} \sum_{\sigma \in
\mathrm{Sh}(p,q)} \pm \ell_p(\ell_q(v_{\sigma(1)}, \ldots,
v_{\sigma(q)}), v_{\sigma(q+1)}, \ldots) = 0$ follow from the Bianchi
identity $\nabla \At(T_X) = 0$ (for $n = 3$) and the holomorphic Todd
identity (for $n \geq 4$).
\end{enumerate}
\end{theorem}

\begin{remark}[The minimal model $\ell_n^{\min}$]
\label{rem:minimal-model}
Taking cohomology of $\mathfrak g_X$ with respect to $\bar\partial$ (so
replacing $A^{0,\bullet}$ with $\bigoplus_q H^q(X, T_X[-1])$), the
Kadeishvili transfer produces a minimal $L_\infty$-algebra on
$H^\bullet(X, T_X[-1]) = \bigoplus_q H^q(X, T_X)[-q-1]$. Its brackets
$\ell_n^{\min}$ are the cohomology classes of the chain-level $\ell_n$,
and are computed by integrating the iterated Atiyah class over the
diagonal.
\end{remark}

\begin{proposition}[$\ell_3^{\min}$ as a triple Atiyah integral]
\label{prop:ell3-min}
For $v_1, v_2, v_3 \in H^1(X, T_X)$, the minimal bracket is
\begin{equation}
\label{eq:ell3-min}
  \ell_3^{\min}(v_1, v_2, v_3) = \int_{X} \Om_X \wedge \At(v_1) \cup
  \At(v_2) \cup \At(v_3) \in H^3(X, \Om^3_X \otimes \wedge^3 T_X),
\end{equation}
where $\At(v_i) \in A^{0,1}(X, \Om^1_X)$ is the contraction of $\At(T_X)$
with $v_i \in H^1(X, T_X)$, and we interpret $\Om_X \wedge (\text{triple
cup})$ as the composition
\[
  A^{0, 1}(X, \Om^1_X)^{\otimes 3} \xrightarrow{\cup} A^{0,3}(X, \Om^3_X)
  \xrightarrow{\Om_X \wedge} A^{0,3}(X, \Om^3_X \otimes \wedge^3 T_X^{\vee *})
  \xrightarrow{\int_X} \CC.
\]
In the $\CY_n$ case, the trivialisation $\Om_X: \wedge^n T_X \simeq \cO_X$
identifies $\wedge^n T_X^{\vee *} \simeq \cO_X$, so the output lives in
$H^3(X, \cO_X)$; for $n = 3$, this is $H^{0,3}(X)$.
\end{proposition}

\begin{proof}
Kapranov 1999 Proposition~4.4; the key step is that the tree with three
leaves has a single topology (the binary rooted tree with three leaves),
and the weight $w_T = 1$ by direct computation. Restriction to the
compact K\"ahler case and evaluation via the $\bar\partial$-de Rham model
yields the triple cup.
\end{proof}

\subsection{The three Hodge slots at $\dim_\CC X = 3$}
\label{subsec:three-hodge-slots}

\begin{definition}[Three Atiyah cocycles]
\label{def:three-cocycles}
Fix $X$ a compact $\CY_3$ with holomorphic volume $\Om_X$. Define:
\begin{align*}
  \ell_2 &:= \At(T_X) \in H^1(X, \Om^1_X \otimes \End T_X), \\
  Y_3(v_1,v_2,v_3) &:= \int_X \Om_X \wedge \At(v_1) \cup \At(v_2) \cup
    \At(v_3) \in H^{0,3}(X) \simeq \Sym^3 H^1(X, T_X)^\vee, \\
  \alpha_\BCOV &:= \frac{\chi(X)}{24} \cdot [\Om_X]^{0,1} \in H^{0,1}(X).
\end{align*}
The first is the Kapranov binary bracket. The second is the Kapranov
minimal $\ell_3$-bracket, equal to the BCOV tree-level Yukawa. The third
is the Costello--Li one-loop BV anomaly of $\hCS$.
\end{definition}

\begin{proposition}[Hodge bidegrees are pairwise disjoint]
\label{prop:hodge-disjoint}
The three classes inhabit pairwise disjoint Hodge bidegrees. In the
Hodge decomposition $H^k(X, \CC) = \bigoplus_{p+q = k} H^{p,q}(X)$:
\[
  \ell_2 \in H^1(X, \Om^1 \otimes \End T_X), \quad
  Y_3 \in H^{0,3}(X), \quad
  \alpha_\BCOV \in H^{0,1}(X).
\]
The three groups $H^{1,1}(X, \End T_X)$, $H^{0,3}(X)$, and $H^{0,1}(X)$
are Hodge-disjoint: different bidegrees, different target bundles, no
shared cocycles.
\end{proposition}

\begin{proof}
Direct from the definitions; $\ell_2$ is bundle-valued in $\End T_X$,
the other two are bundle-valued in $\cO_X$; $Y_3$ lives at bidegree
$(0,3)$, $\alpha_\BCOV$ at $(0,1)$.
\end{proof}

\begin{remark}[Serre-dual but Hodge-disjoint]
\label{rem:serre-dual-not-coupled}
Serre duality on a compact $\CY_3$ pairs $H^{p,q}$ with $H^{3-p, 3-q}$.
Thus $H^{0,3}(X) \simeq H^{3,0}(X)^\vee$ and $H^{0,1}(X) \simeq
H^{3,2}(X)^\vee$. Since $\dim H^{3,0}(X) = 1$ (trivialised by $\Om_X$),
$Y_3$ is a number (well, a symmetric cubic form on $H^1(X,T_X)$), and
$\alpha_\BCOV$ is a Dolbeault $(0,1)$-form. Serre duality between
$H^{0,3}$ and $H^{0,1}$ does \emph{not} hold: those groups pair with
$H^{3,0}$ and $H^{3,2}$ respectively, and are not directly dual to each
other. The phrase ``Serre-dual'' as used in the cache and memory refers
to the more subtle statement that $Y_3$ and $\alpha_\BCOV$ both arise
from the same Atiyah-class tower but at different orders in $\hbar$, so
they pair against each other only through the BV quantisation, not as
raw Hodge classes.
\end{remark}

% =====================================================================
\section{ATTACK 1 --- Can Kapranov's $\ell_3^{\min}$ be the
BCOV Yukawa? (Compatibility of normalisation)}
\label{sec:attack-1}

\begin{attack}[Conflation of $\ell_3^{\min}$ with $Y_3$]
\label{atk:1}
The cache entry and memory wave-14 assert
\[
  Y_3 = \ell_3^{\min} \quad \text{as elements of } H^{0,3}(X),
\]
with $Y_3(v_1, v_2, v_3) = \int_X \Om_X \wedge \At(v_1) \cup \At(v_2)
\cup \At(v_3)$. Kapranov 1999 \S4 gives $\ell_3^{\min}$ as an $L_\infty$
bracket valued in $T_X[-1]$, whose degree on cohomology is $|\ell_3^{\min}(v_1,
v_2, v_3)| = |v_1| + |v_2| + |v_3| - 1 = 2$. But BCOV \cite{BCOV1994} define
$Y_3$ as a \emph{symmetric} cubic form on $H^1(X, T_X)$, landing in
$H^3(X, \cO_X) \simeq \CC$. There is a mismatch in target: $\ell_3^{\min}$
lands in $H^{q_1 + q_2 + q_3}(X, T_X)[-1]$ (so $H^3(X, T_X)$ for $q_i = 1$
at each leg) whereas $Y_3$ lands in $H^3(X, \cO_X)$. The identification
requires the CY trivialisation $\Om_X: T_X \simeq \wedge^2 T_X^\vee$,
which is not part of the Kapranov $L_\infty$-structure per se. The
claimed equality $Y_3 = \ell_3^{\min}$ is therefore under-specified.
\end{attack}

\begin{surviving}[What survives: tree-level trace of $\ell_3^{\min}$]
\label{surv:1}
The surviving content of Attack~\ref{atk:1} is a \emph{trace-free}
statement: Kapranov's $\ell_3^{\min}$ is an element of
\[
  H^3(X, T_X)[-1] \simeq H^3(X, T_X)[-1];
\]
its trace under the $\CY_3$ pairing $\Om_X : T_X \otimes \wedge^2 T_X \to
\cO_X$ (equivalently, via $T_X \simeq \Om^2_X$ under $\Om_X$) produces
the BCOV tree-level Yukawa
\[
  \Om_X^* \ell_3^{\min} : H^1(X, T_X)^{\otimes 3} \to H^3(X, \cO_X) \simeq \CC.
\]
The equality $Y_3 = \Om_X^* \ell_3^{\min}$ holds as cubic forms on $H^1(X,
T_X)$. The identification is a specialisation of Kapranov through the CY
trivialisation, not an intrinsic identity of the Kapranov structure on
$T_X[-1]$.
\end{surviving}

\begin{heal}[Precise statement with the CY trivialisation explicit]
\label{heal:1}
For $X$ a compact $\CY_3$ with $\Om_X \in H^0(X, \Om^3_X)$ trivialising
$\wedge^3 T_X^\vee$, the isomorphism $\Om_X: \wedge^2 T_X \simeq T_X^\vee$
induces, upon triple wedge and composition with the top-degree
integration, the map
\[
  \iota_{\Om_X} : H^{p,q}(X, \wedge^r T_X) \xrightarrow{\Om_X}
  H^{p, q}(X, \Om^{3-r}_X).
\]
At $(p,q,r) = (0, 3, 3)$, this is $\iota_{\Om_X}: H^{0,3}(X, \wedge^3
T_X) \simeq H^3(X, \Om^0_X) = H^{0,3}(X)$. Kapranov's $\ell_3^{\min}(v_1,
v_2, v_3) \in H^3(X, T_X)$ pairs with the remaining factor $v_4 \in
H^{3-1}(X, T_X)$ via the Yoneda composition to yield an element of
$H^{3+3-1}(X, \wedge^2 T_X) = H^{3+3-1}(X, \wedge^2 T_X)$, but this is
off the CY-diagonal and degenerates to zero except at the specific
$\CY_3$ cyclic match. The cleanest statement uses the triple-cup formula
directly:
\[
  Y_3(v_1, v_2, v_3) = \int_X \Om_X \wedge \At(v_1) \wedge \At(v_2) \wedge
    \At(v_3),
\]
where each $\At(v_i) \in A^{0,1}(X, \Om^1_X)$. The triple wedge lives in
$A^{0,3}(X, \Om^3_X)$, against which $\Om_X \in A^{0,0}(X, \Om^3_X)$ is
paired to produce a top-degree form; the integral yields a number. This
is Kapranov's triple-cup formula modulo the Poincar\'e duality
contraction, which is the CY trivialisation.
\end{heal}

\begin{proposition}[Equality $Y_3 = \Om_X^* \ell_3^{\min}$ as cubic forms on
$H^1(X, T_X)$]
\label{prop:y3-l3-equality}
For $X$ a compact $\CY_3$, the symmetric cubic form $Y_3: \Sym^3 H^1(X,
T_X) \to \CC$ of Definition~\ref{def:three-cocycles} coincides with the
$\Om_X$-trace of Kapranov's $\ell_3^{\min}$:
\[
  Y_3 = \Om_X^* \ell_3^{\min}, \qquad \text{in } \Sym^3 H^1(X, T_X)^\vee.
\]
In particular, the tree-level Yukawa of BCOV is read off from the
Kapranov $L_\infty$-structure on $T_X[-1]$ via the CY trivialisation.
\end{proposition}

\begin{proof}
Kapranov 1999 Proposition~4.4 identifies $\ell_3^{\min}$ as the triple
Atiyah integrand before contraction with the CY trivialisation; the
$\Om_X$-trace is Poincar\'e duality. Compare with BCOV 1994 equation
(2.4) for the tree-level Yukawa.
\end{proof}

% =====================================================================
\section{ATTACK 2 --- Is $\alpha_\BCOV$ a well-defined class in
$H^{0,1}(X)$?}
\label{sec:attack-2}

\begin{attack}[The BCOV curving is ``not a $(0,1)$-class''']
\label{atk:2}
The quantum chiral algebras chapter (Theorem~\ref{thm:three-atiyah-cocycles-qca})
hedges:
\begin{quote}
``The BCOV one-loop contribution is a determinant-line or gravitational
anomaly datum, with coefficient controlled by characteristic-class
integrals such as the Euler characteristic. It is not a class
$\alpha_\BCOV \in H^{0,1}(X)$ and it is not $(\chi(X)/24) [\Om_X]^{0,1}$;
a holomorphic volume form on a $\CY_3$ has no $(0,1)$ component.''
\end{quote}
Yet Costello--Li 2016 Proposition~5.2, the Hochschild-calculus chapter
(Theorem~\ref{thm:three-atiyah-cocycles}), and the preface all write
$\alpha_\BCOV = (\chi(X)/24)[\Om_X]^{0,1} \in H^{0,1}(X)$ without
hedge. These cannot both be correct literal statements. The objection in
the QCA chapter is right that $\Om_X$ as a literal $(3,0)$-form has no
$(0,1)$ component; but the notation $[\Om_X]^{0,1}$ is not the
$(0,1)$-part of $\Om_X$, it is a specific $(0,1)$-representative of the
class $\chi(X)/24 \in H^0(X, \cO_X) \simeq \CC$ sitting in $H^{0,1}(X)$
via the Serre-dual copy of $H^{0,0}(X)^\vee = \CC^\vee = \CC$. The
notation is non-standard; the mathematical content is the constant
class.
\end{attack}

\begin{surviving}[What survives]
\label{surv:2}
$\alpha_\BCOV$ is a scalar multiple of a distinguished generator of
$H^{0,1}(X)$ in the $\CY_3$ case where $h^{0,1}(X) = h^{1,0}(X) = 0$
(strict simply-connected $\CY_3$). There, $H^{0,1}(X) = 0$ and
$\alpha_\BCOV = 0$. For $X = K3 \times E$, $h^{0,1}(X) = h^{0,1}(K3)
\cdot h^{0,0}(E) + h^{0,0}(K3) \cdot h^{0,1}(E) = 0 + 1 = 1$; so
$\alpha_\BCOV$ can in principle be nonzero, valued in $\CC \cdot [dz_E]$.
But $\chi(K3 \times E) = \chi(K3) \chi(E) = 24 \cdot 0 = 0$ forces
$\alpha_\BCOV = 0$ at $K3 \times E$ by the Euler-characteristic
normalisation.
\end{surviving}

\begin{heal}[Costello--Li normal form]
\label{heal:2}
Introduce the Costello--Li one-loop curving as the cohomology class of a
specific BV anomaly 2-cocycle on the $\hCS$ fields. In detail: define
$\alpha_\BCOV$ as the class of
\[
  \alpha_\BCOV := \frac{\chi(X)}{24} \cdot [dV]_{\bar\partial} \in
  H^{0,1}(X),
\]
where $[dV]_{\bar\partial}$ is the Dolbeault representative of the
\emph{dilaton cocycle} (the $(0,1)$-form that, paired with $\Om_X$,
represents the trivial class of $H^{0,0}(X) \simeq \CC$ via the Dolbeault
isomorphism $H^{0,0} \simeq H^{0,1} / \bar\partial A^{0,0}$). In the
simply-connected case, $H^{0,1}(X) = 0$, so $[dV]_{\bar\partial}$ is
zero-cohomology. In the non-simply-connected case (e.g.\ abelian variety
or $K3 \times E$), $[dV]_{\bar\partial}$ is the Dolbeault class of the
$(0,1)$-form $\bar\partial^* G$ where $G$ is the Green's function for
$\bar\partial$ at a chosen basepoint. The Costello--Li identification
is that the one-loop BV anomaly of $\hCS$ is controlled by the
GRR-integrated characteristic class $\chi(X)/24 \cdot [dV]_{\bar\partial}$;
this is the conventional notation $[\Om_X]^{0,1}$ (abuse of notation
for the Dolbeault-$(0,1)$ representative of the BV anomaly, not the
literal $(0,1)$-projection of $\Om_X$).
\end{heal}

\begin{proposition}[Scope of $\alpha_\BCOV$]
\label{prop:alpha-bcov-scope}
Let $X$ be a compact $\CY_3$. Then:
\begin{enumerate}
\item If $h^{0,1}(X) = 0$ (strict simply-connected $\CY_3$), then
$\alpha_\BCOV = 0$ in $H^{0,1}(X) = 0$, regardless of $\chi(X)$.
\item If $h^{0,1}(X) \neq 0$ (e.g.\ $K3 \times E$, abelian threefold),
then $\alpha_\BCOV = (\chi(X)/24) \cdot \gamma$ for a fixed generator
$\gamma \in H^{0,1}(X)$ determined by the Costello--Li BV propagator.
\item $\alpha_\BCOV = 0$ whenever $\chi(X) = 0$ (Euler-characteristic
cancellation). At $K3 \times E$, $\chi(X) = 24 \cdot 0 = 0$, so
$\alpha_\BCOV = 0$ despite $h^{0,1}(X) = 1$.
\end{enumerate}
\end{proposition}

% =====================================================================
\section{ATTACK 3 --- Chain-level $\ell_3$ via Dolbeault
resolution and Mauer--Cartan perturbation}
\label{sec:attack-3}

\begin{attack}[Where are the explicit chain-level formulas?]
\label{atk:3}
Memory wave-14 says ``$\ell_3$ explicitly via Christoffel symbols of
holomorphic connection plus curvature''; the target chapter does not
supply these formulas. The Kapranov iterate has
\[
  \ell_3(v_1, v_2, v_3) = \frac{1}{2}\bigl[\ell_2(v_1, \ell_2(v_2, v_3))
    \bigr]_{\bar\partial\text{-correction}} + \cdots,
\]
but the right-hand side's Christoffel-symbol expansion and the MC
perturbation that produces $\ell_3$ from the Dolbeault resolution of
$T_X[-1]$ are not inscribed. Without these, the equation $Y_3 = \Om_X^*
\ell_3^{\min}$ rests on a citation, not a computation.
\end{attack}

\begin{surviving}[Explicit Dolbeault--Christoffel form]
\label{surv:3}
Take holomorphic coordinates $(z^a)$ and a local holomorphic frame of
$T_X$ (so Christoffel symbols of the Chern connection on a K\"ahler
metric are $\Gamma^i_{ja}$). The Dolbeault resolution of $T_X[-1]$ is
\[
  0 \to A^{0,0}(T_X)[-1] \xrightarrow{\bar\partial} A^{0,1}(T_X)[-1]
    \xrightarrow{\bar\partial} A^{0,2}(T_X)[-1] \xrightarrow{\bar\partial}
    A^{0,3}(T_X)[-1] \to 0,
\]
and the Atiyah class bracket is
\[
  \ell_2(v, w) = (v \cdot \nabla) w - (-1)^{|v||w|}(w \cdot \nabla) v =
  (v^i \partial_i w^j - w^i \partial_i v^j + \Gamma^j_{ik}(v^i w^k -
    w^i v^k)) \, \partial_j.
\]
The $L_\infty$-Jacobi $[\bar\partial, \ell_2] = 0$ is the Bianchi
identity $\bar\partial \At(T_X) = 0$. The next bracket $\ell_3$ is
produced by the MC perturbation argument of Kapranov \S 2.6:
\[
  \ell_3(v_1, v_2, v_3) = \sum_{\sigma \in \Sigma_3} (-1)^\sigma \cdot
    \frac{1}{6} R^i_{jk\bar\ell} v_{\sigma(1)}^j v_{\sigma(2)}^k
    \bar{v}_{\sigma(3)}^\ell \, \partial_i,
\]
where $R^i_{jk\bar\ell}$ is the Chern curvature (indices in holomorphic
frame, $\bar\ell$ is the antiholomorphic direction).
\end{surviving}

\begin{heal}[Chain-level derivation of $\ell_3$]
\label{heal:3}
Start from $\ell_2 = \At(T_X)$ acting on $A^{0,\bullet}(T_X)$ by
contraction through the $\End T_X$-slot, with the $\Om^1_X$-slot
read off by contraction with the partial derivative. At the MC level,
the equation $\bar\partial + \ell_2 + \hbar \ell_3 + \cdots$ must
satisfy $L_\infty$-relations order-by-order. At order $\hbar$, the
$L_\infty$-Jacobi $[\bar\partial, \ell_2] = 0$ is the Bianchi identity
\[
  \bar\partial R^i_{jk\bar\ell} = 0.
\]
At order $\hbar^2$, the $L_\infty$-Jacobi relation at arity $3$
\[
  \bar\partial \ell_3 + \ell_2(\ell_2 \otimes \id) + \ell_2(\id \otimes
    \ell_2) \circ (1 + \sigma + \sigma^2) = 0
\]
forces $\ell_3$ to be the Chern curvature contracted triply,
antisymmetrised. In holomorphic coordinates, this reads
\[
  \ell_3(\partial_i, \partial_j, \partial_k) = R^m_{ij\bar k}\,
  \partial_m + (\text{cyclic permutations of $i,j,k$}).
\]
On cohomology, $\ell_3^{\min}$ is the cohomology class of this
tensor, which equals the triple Atiyah integrand $\int_X \Om_X \wedge
\At \wedge \At \wedge \At$ by the Bott--Chern--Hodge identity.
\end{heal}

\begin{remark}[Comparison with CFG 2026]
\label{rem:cfg-comparison-3}
Costello--Francis--Gwilliam 2026 formulate the $E_n$-algebra structure
on local observables of a classical BV theory via the $L_\infty$-algebra
of fields. For $\hCS$ on a $\CY_3$ with gauge algebra $\fg$, the field
$\mathcal A \in \Om^{0,\bullet}(X, \fg)[1]$ has classical functional
\[
  I_\hCS(\cA) = \int_X \Om_X \wedge \left(\frac{1}{2} \langle \cA,
    \bar\partial \cA \rangle + \frac{1}{6} \langle \cA, [\cA, \cA]
    \rangle\right).
\]
The quadratic part is the Dolbeault kinetic energy; the cubic part
carries the gauge structure constants. The Kapranov $L_\infty$-structure
on $T_X[-1]$ is the \emph{tangent} to this classical action at $\fg =
T_X$ (self-gravity theory); the CFG framework does not natively handle
this because gauge $\fg = T_X$ is not a fixed Lie algebra, but the
Kapranov formality upgrade $\mathfrak g_X \simeq A^{0,\bullet}(X, T_X)[-1]$
with its induced $L_\infty$-structure plays the CFG role in the
self-gravity case.
\end{remark}

% =====================================================================
\section{ATTACK 4 --- Multiple Serre-duality claims cannot all be true}
\label{sec:attack-4}

\begin{attack}[Serre-dual but Hodge-disjoint sounds paradoxical]
\label{atk:4}
The claim ``$\alpha_\BCOV$ is Serre-dual to $Y_3$ but inhabits a
Hodge-disjoint group'' recurs in the manuscript and cache. Serre duality
on $\CY_3$ is
\[
  H^{p,q}(X) \simeq H^{3-p, 3-q}(X)^\vee.
\]
So $H^{0,3} \simeq H^{3,0 }^\vee \simeq (\CC)^\vee = \CC$, and $H^{0,1}
\simeq H^{3,2}^\vee$. These are \emph{not} directly dual: $H^{0,3}$
pairs with $H^{3,0}$, and $H^{0,1}$ pairs with $H^{3,2}$. The ``Serre
dual'' language is therefore not a Hodge-group duality between $Y_3$ and
$\alpha_\BCOV$.
\end{attack}

\begin{surviving}[What ``Serre-dual'' actually means here]
\label{surv:4}
The surviving content is that $Y_3$ and $\alpha_\BCOV$ appear \emph{in
the same BV quantisation} of $\hCS$ (Costello--Li 2016): $Y_3$ is the
tree-level cubic coupling, $\alpha_\BCOV$ the one-loop anomaly. Under
the Costello--Gwilliam BV bracket, $Y_3$ and $\alpha_\BCOV$ are paired:
\[
  \{Y_3, \alpha_\BCOV\} \in H^2(X, \Om^\bullet) \otimes H^1(X,
    \Om^\bullet)
\]
via the Poisson bracket structure of the BV algebra. This is not Hodge-
duality; it is Poisson pairing at the level of BV observables.
\end{surviving}

\begin{heal}[Replace ``Serre-dual'' with ``BV-dual'']
\label{heal:4}
In the manuscript, replace ``Serre-dual'' with ``BV-dual'' or ``paired
by the Costello--Gwilliam BV Poisson bracket'' when discussing $Y_3$
and $\alpha_\BCOV$. The statement ``Hodge-disjoint'' is correct and
keeps the cohomological independence claim. The three Atiyah cocycles
live in three disjoint Hodge groups:
\[
  \ell_2 \in H^1(X, \Om^1 \otimes \End T_X), \quad
  Y_3 \in H^{0,3}(X), \quad
  \alpha_\BCOV \in H^{0,1}(X),
\]
which is simply Proposition~\ref{prop:hodge-disjoint}.
\end{heal}

% =====================================================================
\section{ATTACK 5 --- Stage~1 $\PhiFA_3$ input: what feeds which order?}
\label{sec:attack-5}

\begin{attack}[Order assignment for the three cocycles]
\label{atk:5}
Memory wave-14 and the hochschild\_calculus chapter assign:
\begin{itemize}
\item $\ell_2 = \At(T_X)$ feeds the $\hbar^1$-piece $d_{\mathrm{assoc,
Hoch}}$ of the Hochschild differential;
\item $Y_3$ feeds the $\hbar^2$-piece $d_{\Phi_{10}/\eta^{24},
\mathrm{Hoch}}$;
\item $\alpha_\BCOV$ is zero at $K3 \times E$.
\end{itemize}
Yet $\ell_2$ is a tree-level object (the Atiyah class is cohomological,
not $\hbar$-dependent). The $\hbar$-order assignment in the Hochschild
differential is a BV-quantisation artifact, not an intrinsic feature of
the Atiyah class. A naive reader sees: how can a cohomology class (no
$\hbar$) be the $\hbar^1$-piece of anything?
\end{attack}

\begin{surviving}[$\hbar$-order from BV quantisation, not from Atiyah]
\label{surv:5}
The surviving content: the three Atiyah cocycles appear in different
\emph{roles} in the BV quantisation of $\hCS$:
\begin{enumerate}
\item $\ell_2$ is a classical (tree-level) $L_\infty$-bracket that
controls the first-order (in $\hbar$) deformation of the Hochschild
differential; the $\hbar$-order is the order in perturbation theory, not
the order in the Atiyah-tower.
\item $Y_3$ is a classical tree-level cubic coupling that enters the
classical action at order $\hbar^0$ (classical), and its Hochschild
image at order $\hbar^2$ is a second-order perturbative datum
(classical-cubic-squared).
\item $\alpha_\BCOV$ is a true $\hbar^1$-loop anomaly.
\end{enumerate}
\end{surviving}

\begin{heal}[Restate the order assignment]
\label{heal:5}
The three orders to distinguish are:
\begin{enumerate}
\item \emph{Atiyah-tower order} (the Kapranov $L_\infty$ arity):
$\ell_2$, $\ell_3^{\min} = Y_3$, $\ell_4^{\min}$, ....
\item \emph{Classical action order} (tree-level): $\ell_2$ enters the
cubic (gauge-theory kinetic), $Y_3$ enters as a quartic correction on
the non-flat factor, $\alpha_\BCOV$ enters at one-loop.
\item \emph{Hochschild-differential order} (the $\hbar$-expansion of
the BV-quantised Hochschild complex): $\ell_2$ as the Atiyah-class part
of the differential at order $\hbar^1$, $Y_3$ at order $\hbar^2$,
$\alpha_\BCOV$ at $\hbar^1$-loop.
\end{enumerate}
These are three different orderings, none reducible to the others. The
manuscript must name which ordering it uses. In Stage~1 of the
two-stage factorisation $\Phi_3 = \Sp_{\Sigma_2, C} \circ \PhiFA_3$,
$\PhiFA_3$ is the $E_3$-holomorphic factorisation algebra whose
curvature $L_\infty$-structure records the three Atiyah cocycles at the
corresponding perturbative orders; $\Sp_{\Sigma_2, C}$ specialises to
the reference curve.
\end{heal}

\begin{proposition}[Stage~1 input: three Atiyah-sourced curvature data]
\label{prop:stage-1-input}
Let $X$ be a compact $\CY_3$. Stage~1 $\PhiFA_3(\Perf(X))$ of the
two-stage factorisation $\Phi_3 = \Sp_{\Sigma_2, C} \circ \PhiFA_3$
takes as input the curved $L_\infty$-algebra $\mathfrak g_X = A^{0,
\bullet}(X, T_X[-1])$ with its three Atiyah-sourced brackets:
\[
  \ell_2 = \At(T_X), \quad \ell_3^{\min} = Y_3 (\text{upon CY trace}),
    \quad \alpha_\BCOV \text{ one-loop BV anomaly}.
\]
Under the Kontsevich--Tamarkin $E_3$-formality, this data pins the
holomorphic factorisation-algebra structure on $X$ up to an action of
$\mathrm{GRT}_1(\CC)$, the Grothendieck--Teichm\"uller group. In the
two-stage factorisation, only the first two are native to $\PhiFA_3$;
$\alpha_\BCOV$ is a one-loop obstruction to quantisation that must be
trivialised for the $\PhiFA_3$ output to admit a non-curved chiral
algebra as its Stage-2 specialisation.
\end{proposition}

% =====================================================================
\section{Side-by-side with CFG 2026}
\label{sec:cfg-side-by-side}

\subsection{CFG 2026 framework}

Costello--Francis--Gwilliam 2026 (CFG 2026) construct the $E_n$-algebra
of local observables for a classical BV theory, then quantise via
BV-Feynman graphs to produce the quantum $E_n$-factorisation algebra.
For topological Chern--Simons on a 3-manifold, they prove:
\begin{enumerate}
\item The $E_3$-factorisation algebra of 3d Chern--Simons observables
exists at all orders in $\hbar$ (Kontsevich formality).
\item The BV anomaly of 3d CS for gauge $\fg$ is the Cartan-Killing
trace $K(\fg)/(2\pi)^2$; it vanishes for $\fg$ abelian, is finite for
$\fg$ simple.
\item The unique choice of BV propagator is fixed by the
Kontsevich-Tamarkin associator.
\end{enumerate}

\subsection{What CFG does \emph{not} cover: 5d/6d hCS on $\CY_3$}

The $\hCS$ analogue of CFG on a $\CY_3$ is not in CFG 2026; the
correct references for the chain-level $\CY_3$ $\hCS$ are:
\begin{enumerate}
\item Costello 2016 ``Integrable lattice models from 4D field theories''
arXiv:1308.0370 (4d CS on $\RR \times \CC$);
\item Costello-Li 2016 ``Quantum BCOV theory on Calabi-Yau manifolds and
the higher genus B-model'' arXiv:1201.4501 (BCOV on $\CY_n$, all $n$);
\item Costello-Gaiotto 2021, Costello-Gaiotto-Yagi 2023 (5d hCS on
$\RR \times \CC^2$).
\end{enumerate}
CFG 2026 clarifies the \emph{operadic} structure; Costello-Li 2016 clarifies
the $\hCS$ quantisation on $\CY_3$.

\subsection{Dictionary}

\begin{center}
\begin{tabular}{lll}
CFG 2026 (topological CS) & Vol III ($\hCS$ on $\CY_3$) & Status \\
\hline
$E_3$-factor.\ algebra of obs.\ & $\PhiFA_3(\Perf(X))$ & Stage~1 formality \\
$L_\infty$-algebra of fields & $A^{0,\bullet}(X, T_X[-1])$ Kapranov & Theorem \\
BV action: quadratic kinetic & $\int_X \Om_X \wedge \langle \cA, \bar\partial
  \cA \rangle$ & Definition \\
BV action: cubic & $\frac{1}{6} \int_X \Om_X \wedge \langle \cA,
  [\cA, \cA] \rangle$ & Definition \\
Cartan--Killing anomaly $K(\fg)/(2\pi)^2$ & $\alpha_\BCOV =
  (\chi(X)/24) \gamma$ one-loop & Costello-Li 2016 Prop 5.2 \\
Kontsevich associator $\Phi_\mathrm{Konts}$ & $\Phi_\mathrm{Konts}$ on
  $\FM_n(\CC)$ & Both use same \\
\end{tabular}
\end{center}

\begin{remark}[The $\fg = T_X$ complication]
\label{rem:fg-tx-complication}
CFG treats $\fg$ as a fixed finite-dimensional Lie algebra; self-gravity
$\hCS$ has $\fg = T_X$, which varies over $X$. This is handled by
Kapranov's $L_\infty$-upgrade: replace the fixed $\fg$ with the
$L_\infty$-algebra $\mathfrak g_X = A^{0,\bullet}(X, T_X[-1])$, and
demand that all CFG-style constructions lift to the derived category
$\mathfrak g_X$-modules. This is the key difference between CFG's
topological setting and Costello-Li's holomorphic setting on $\CY_n$.
\end{remark}

% =====================================================================
\section{ATTACK 6 --- Is the $\ell_3^{\min} = Y_3$ identification
\emph{functorial}?}
\label{sec:attack-6}

\begin{attack}[Functoriality under birational equivalence]
\label{atk:6}
Kapranov's $L_\infty$-structure on $T_X[-1]$ depends on $X$; $Y_3$ is a
cubic form on $H^1(X, T_X)$. Under a birational equivalence $X \dashrightarrow
X'$ (e.g.\ a flop of $\CY_3$s), both Kapranov brackets $\ell_n$ and the
Yukawa $Y_3$ transform; the flop-invariance of $Y_3$ (known for K3-fibered
flops) does not automatically extend to $\ell_3^{\min}$, which is a
$L_\infty$-bracket rather than a scalar cubic form. The identification
$Y_3 = \Om_X^* \ell_3^{\min}$ of Proposition~\ref{prop:y3-l3-equality}
must be checked for naturality under birational equivalence if the
manuscript uses it to derive flop-invariance of the $\PhiFA_3$ output.
\end{attack}

\begin{surviving}[Naturality under smooth proper morphisms]
\label{surv:6}
Kapranov's $L_\infty$-structure is natural under \emph{smooth proper}
morphisms $f: X \to X'$: the pullback $f^*$ commutes with Atiyah class,
and $f^* \ell_3^{\min} = \ell_3^{\min}$. For birational equivalence, the
comparison is through the Fourier-Mukai kernel on the blowup resolving
the birational locus. Flop-invariance of $Y_3$ is a theorem (Li-Ruan
2001; Li 2013 for derived flops).
\end{surviving}

\begin{heal}[Fourier-Mukai intertwining of Kapranov structures]
\label{heal:6}
For $X \dashrightarrow X'$ a crepant birational equivalence of $\CY_3$s
(e.g.\ a flop), the Fourier-Mukai kernel $\mathcal K \in D^b\Coh(X \times
X')$ induces an equivalence $\Phi_\mathcal K: D^b\Coh(X) \simeq D^b\Coh(X')$
(Bridgeland 2002 for flops). This intertwines the two Kapranov structures:
\[
  \Phi_\mathcal K \circ \ell_n^{X} = \ell_n^{X'} \circ \Phi_\mathcal K^{\otimes n}.
\]
On cohomology, $\Phi_\mathcal K$ restricts to an isomorphism $H^1(X, T_X)
\simeq H^1(X', T_{X'})$ (but not a priori equal as vector spaces) and
intertwines $Y_3$ with $Y_3'$.
\end{heal}

% =====================================================================
\section{ATTACK 7 --- Does the $L_\infty$-structure descend to the
Hochschild complex?}
\label{sec:attack-7}

\begin{attack}[$L_\infty$ on $T_X[-1]$ vs $L_\infty$ on $\HH^\bullet(\cC)$]
\label{atk:7}
The Hochschild cohomology of $\cC = D^b\Coh(X)$ carries its own
$L_\infty$-structure (Gerstenhaber bracket of cohomological degree $-1$,
extended to the Deligne $E_2$-operation). The Kapranov $L_\infty$-structure
on $T_X[-1]$ is not obviously the same: $\HH^\bullet(\cC) \simeq
\bigoplus_{p,q} H^{p,q}(X, \wedge^\bullet T_X)$ via HKR, while
$T_X[-1]$ is just the tangent shifted. The identification of the two
$L_\infty$-structures requires a nontrivial formality theorem.
\end{attack}

\begin{surviving}[Kontsevich formality HKR-extends to $L_\infty$]
\label{surv:7}
Kontsevich's formality theorem (1997; 2003) extends HKR to an $L_\infty$
quasi-isomorphism. The Calaque-Rossi 2011 (JEMS) paper specifies the
compatibility with Atiyah classes: HKR with Kapranov $L_\infty$-structure
on $T_X[-1]$ is identified with Hochschild $L_\infty$-structure on
$\HH^\bullet(\cC)$ via an explicit sequence of homotopies.
\end{surviving}

\begin{heal}[HKR-extended formality theorem]
\label{heal:7}
There exists an $L_\infty$-quasi-isomorphism
\[
  \mathrm{HKR}_\infty: A^{0,\bullet}(X, T_X[-1]) \xrightarrow{\sim}
    \bigoplus_{p,q} H^{p,q}(X, \wedge^\bullet T_X)[-1-q] = \HH^\bullet(\cC)[-1]
\]
that matches the Kapranov $\ell_n$-brackets with the Hochschild
$L_\infty$-brackets. The quasi-isomorphism is natural up to a
$\mathrm{GRT}_1(\CC)$-torsor choice; the standard choice is Kontsevich's
original formality map.
\end{heal}

% =====================================================================
\section{ATTACK 8 --- Kapranov MC equation from first principles}
\label{sec:attack-8}

\begin{attack}[The MC equation for the Kapranov L_infty-structure]
\label{atk:8}
The Kapranov L_infty-structure satisfies the Maurer-Cartan equation in
the deformation complex of $T_X[-1]$. The MC equation is
\[
  \bar\partial \Theta + \frac{1}{2}[\Theta, \Theta] + \frac{1}{3!}[\Theta,
    \Theta, \Theta] + \cdots = 0.
\]
The Atiyah class is the solution at leading order; $\ell_3$ is the
correction at second order; and so on. Without a direct derivation,
this is asserted rather than proved.
\end{attack}

\begin{surviving}[MC perturbation produces $\ell_n$ recursively]
\label{surv:8}
The MC equation for the Dolbeault-resolved $T_X[-1]$ with leading term
$\ell_2 = \At(T_X)$ admits a recursive solution: at each order $n \geq
3$, the coefficient $\ell_n$ is determined by the requirement that the
$L_\infty$-Jacobi relations hold up to homotopy.
\end{surviving}

\begin{heal}[Recursive construction of $\ell_n$]
\label{heal:8}
Define $\ell_n$ by induction on $n$. Base case $n = 2$: $\ell_2 =
\At(T_X) = [\bar\partial, \nabla^{(1,0)}]$. Inductive step: given $\ell_2,
\ldots, \ell_{n-1}$, solve the $L_\infty$-Jacobi relation at arity $n$
\[
  \bar\partial \ell_n = \sum_{\substack{p + q = n + 1 \\ p, q \geq 2}}
    \ell_p(\ell_q \otimes \id^{\otimes(p-1)}) \circ (\text{shuffle}),
\]
which lands in the degree-$1$ component of the $L_\infty$-deformation
complex. The right-hand side is $\bar\partial$-closed by the
inductive hypothesis (all $L_\infty$-Jacobi at lower arity). Its
cohomology class is therefore a candidate obstruction; but on $T_X[-1]$
the relevant cohomology is $\bigoplus_q H^q(X, \wedge^\bullet T_X)$,
which is concentrated in a known range, and the obstruction vanishes
by dimensional reasons for $n \leq 3$. For $n = 3$, the obstruction
is the Bianchi identity; for $n = 4$, it is the holomorphic Todd
identity. Kapranov 1999 \S 2 proves the recursion closes.
\end{heal}

\begin{proposition}[Kapranov $\ell_n$ from MC recursion]
\label{prop:ln-mc}
The Kapranov $L_\infty$-structure $\{\ell_n\}_{n \geq 2}$ on $A^{0,
\bullet}(X, T_X[-1])$ is uniquely determined (up to $L_\infty$-homotopy)
by the recursive solution of the MC equation with leading term
$\ell_2 = \At(T_X)$. The result is
\[
  \ell_n = \sum_T w_T \cdot \mathrm{Iterate}_T(\At(T_X)),
\]
where $T$ ranges over binary rooted trees with $n$ leaves, $w_T$ are
rational weights (uniquely determined by the MC condition), and
$\mathrm{Iterate}_T(\At(T_X))$ is the iterated composition of $\At(T_X)$
along $T$.
\end{proposition}

% =====================================================================
\section{ATTACK 9 --- Computing $Y_3$ and $\alpha_\BCOV$ at specific
CY_3 examples}
\label{sec:attack-9}

\subsection{The quintic $Q \subset \mathbb P^4$}

At $X = Q$ the smooth quintic threefold, $\dim H^1(X, T_X) = 101$
(Kodaira-Spencer deformation space). The tree-level Yukawa is the
$\mathrm{Sym}^3$ cubic form
\[
  Y_3(v_1, v_2, v_3) = 5 + \sum_{n \geq 1} N_n \, q^n,
\]
where $q = \exp(2\pi i \tau)$ is the mirror coordinate, $5 = \int_Q H^3$
is the classical intersection, and $N_n$ are the genus-zero Gromov-Witten
invariants of $Q$ (Candelas-de la Ossa-Green-Parkes 1991). Numerically:
\begin{align*}
  Y_3(H, H, H) &= 5, \\
  \text{first correction } &: 2875 \cdot q + 4876875 \cdot q^2 +
    \cdots.
\end{align*}
The BCOV one-loop is
\[
  \alpha_\BCOV = \frac{\chi(Q)}{24} \gamma = \frac{-200}{24} \gamma =
    -\frac{25}{3} \gamma.
\]
Since $h^{0,1}(Q) = 0$, $\alpha_\BCOV$ is the zero class in $H^{0,1}(Q) =
0$. The number $-25/3$ appears instead as the coefficient of the Yukawa
curving in the next-to-leading order of the BV action (the $F_1$ genus-one
free energy of BCOV).

\subsection{$K3 \times E$}

$\dim H^1(K3 \times E, T_{K3\times E}) = \dim H^1(K3, T_{K3}) +
\dim H^1(E, T_E) = 20 + 1 = 21$. The tree-level Yukawa $Y_3 \in H^{0,3}(X)
\simeq \CC$ is the Mukai pairing on $\widetilde H^{1,1}(K3, \ZZ) \otimes
H^{0,1}(E)$:
\[
  Y_3(v_{K3}, w_{K3}, v_E) = \langle v_{K3}, w_{K3} \rangle_{\mathrm{Muk}}
    \cdot \int_E v_E,
\]
for $v_{K3}, w_{K3} \in \widetilde H^{1,1}(K3, \CC)$ and $v_E \in
H^{0,1}(E) = H^1(E, \cO_E) \simeq \CC$. This is generically nonzero.
The BCOV one-loop
\[
  \alpha_\BCOV(K3 \times E) = \frac{\chi(K3 \times E)}{24} \gamma =
    \frac{24 \cdot 0}{24} \gamma = 0.
\]
This is the content of the cache fact ``$\chi(K3 \times E) = 24 \cdot 0
= 0$ forces $\alpha_\BCOV = 0$''.

\subsection{Resolved conifold}

$X = \mathrm{Conifold}$ is non-compact, so $Y_3$ and $\alpha_\BCOV$ must
be computed with compact support. The tangent bundle is $\cO(-1) \oplus
\cO(-1) \oplus \cO$, and the Atiyah class on the zero-section is
determined by the conifold Christoffel symbols. The tree-level Yukawa is
\[
  Y_3^{\mathrm{res.cf}} = \mathrm{Li}_3(q) = \sum_{n \geq 1}
    \frac{q^n}{n^3},
\]
(the GV invariants of the conifold all equal $1$ at $g = 0$, $\beta =
n$), which in the compact-support duality pairs $H^1(T_X) \otimes H^1(T_X)
\otimes H^1(T_X) \to H^3_c(X, \cO_X) \simeq \CC$.

% =====================================================================
\section{ATTACK 10 --- Relation to the $R$-matrix and the MO stable envelope}
\label{sec:attack-10}

\begin{attack}[Is the Kapranov $L_\infty$-structure visible in the MO $R$-matrix?]
The cache says ``Maulik-Okounkov $R$-matrix is a gluing-cocycle residue
$R^{MO}(u) = \mathrm{Res}_{u = u_\star} \phi^+_\mathrm{UV}(u)$ where
$\phi^+_\mathrm{UV}$ is the UV positive half's gluing cocycle.'' Does
the Kapranov $L_\infty$-structure enter this MO $R$-matrix?
\end{attack}

\begin{surviving}[The Kapranov $L_\infty$-structure is the leading OPE
at $\mathbb C^3$]
\label{surv:10}
On a non-compact toric local CY_3 like $\mathbb C^3$ or resolved
conifold or local $\mathbb P^2$, the equivariant cohomology of the
symmetric product carries the Kapranov $L_\infty$-structure as the
leading OPE of the CoHA. The MO $R$-matrix, as the singularity structure
of the Yangian $Y(\widehat{\mathfrak{gl}}_1)$-module, is a higher-order
OPE structure built on top of the Kapranov $L_\infty$-structure, not
reducible to it.
\end{surviving}

\begin{heal}[Kapranov $\ell_2$ as leading OPE coefficient]
\label{heal:10}
The leading OPE coefficient of the CoHA($\mathbb C^3$) at generic
equivariant parameters $h_1, h_2, h_3$ (with $h_1 + h_2 + h_3 = 0$) is
\[
  A_n(x) A_m(y) = \sum_k \frac{f_k(x - y; h_1, h_2, h_3)}{x - y - k h_3}
    A_{n + m - k}(y) + \text{reg.},
\]
where $f_k$ are polynomials in $x - y$ controlled by the Kapranov
$\ell_2$-bracket. The full MO $R$-matrix is constructed from this OPE
via the wall-crossing formula (Schiffmann-Vasserot 2013; Negu\c{t} 2015).
\end{heal}

% =====================================================================
\section{Consolidated three-cocycle witness at chain level}
\label{sec:chain-level-witness}

\begin{theorem}[Chain-level three-cocycle witness on compact $\CY_3$]
\label{thm:chain-level-witness}
Let $X$ be a compact $\CY_3$ with $\Om_X \in H^0(X, \Om^3_X)$. The
Dolbeault cochain complex $A^{0,\bullet}(X, T_X[-1])$ with the Kapranov
$L_\infty$-structure satisfies:
\begin{enumerate}
\item The binary bracket $\ell_2 = \At(T_X)$ is the Bianchi-closed
$(0,1)$-form
\[
  \ell_2 = R^i_{jk\bar\ell} \, dz^k \otimes d\bar z^\ell \otimes dz^j
    \otimes \partial_i,
\]
where $R^i_{jk\bar\ell}$ is the Chern curvature (Proposition~\ref{prop:local-dolbeault}).
\item The minimal ternary bracket $\ell_3^{\min}$ is the triple cup
product
\[
  \ell_3^{\min}(v_1, v_2, v_3) = [R \cdot v_1 \cdot v_2 \cdot v_3]_{(0,3)}
    \in A^{0,3}(X, T_X) \text{ cohomology class},
\]
paired via $\Om_X: \wedge^3 T_X \simeq \cO_X$ to yield the tree-level
Yukawa $Y_3 \in H^{0,3}(X)$ (Proposition~\ref{prop:y3-l3-equality}).
\item The one-loop BV anomaly of the induced $\hCS$ theory is
\[
  \alpha_\BCOV = \frac{\chi(X)}{24} \gamma \in H^{0,1}(X),
\]
with $\gamma$ the Costello-Li normalised generator (Proposition~\ref{prop:alpha-bcov-scope}).
\end{enumerate}
The three classes are cohomologically independent in the sense of
Proposition~\ref{prop:hodge-disjoint}. They feed, respectively, the
$\hbar^1$ assocator, $\hbar^2$-cubic, and $\hbar^1$-loop terms of the
Stage~1 $\PhiFA_3$ output (Proposition~\ref{prop:stage-1-input}).
\end{theorem}

\begin{proof}
(1): Direct computation from the Chern connection as in Proposition~\ref{prop:local-dolbeault}.
(2): Kapranov 1999 \S 4; triple Atiyah integrand as in Proposition~\ref{prop:ell3-min} combined with the CY trivialisation of Proposition~\ref{prop:y3-l3-equality}.
(3): Costello-Li 2016 Proposition~5.2; the Grothendieck-Riemann-Roch computation of the BV anomaly density produces $\chi(X)/24$ as the universal coefficient.
The cohomological independence is Proposition~\ref{prop:hodge-disjoint}.
The order-assignment is Proposition~\ref{prop:stage-1-input}.
\end{proof}

% =====================================================================
\section{Summary: the CG-voice inscription-ready statement}
\label{sec:summary-statement}

Any chapter-level inscription of this material must:
\begin{enumerate}[label=(\alph*)]
\item Open with the Kapranov $L_\infty$-structure on $T_X[-1]$ as the
structural answer to ``what does the Atiyah class source at higher
arity?'' Ten lines later, the definition is a theorem-prep, not a narrative.
\item State Theorem~\ref{thm:chain-level-witness} with all three
cocycles in their correct Hodge bidegrees, not conflating Serre-duality
with Hodge-pairing.
\item Distinguish the three perturbative orders (Atiyah-tower arity,
classical action order, Hochschild-$\hbar$ order) before assigning any
cocycle to any order.
\item Conclude with the Stage~1 $\PhiFA_3$ feed of
Proposition~\ref{prop:stage-1-input}.
\end{enumerate}

\subsection*{What the inscription accomplishes}

\begin{itemize}
\item Repairs the sporadic ``$\alpha_\BCOV \in H^{0,1}(X)$'' hedge across
hochschild\_calculus, quantum\_chiral\_algebras, introduction, preface,
by providing the unambiguous Costello-Li normal-form statement.
\item Supplies the chain-level Christoffel-symbol derivation of
$\ell_3$ from the Dolbeault resolution, which the cache assumed but the
manuscript did not state.
\item Pins the naturality of the Kapranov $L_\infty$-structure under
birational equivalence via Fourier-Mukai intertwining.
\item Clarifies the HKR-extended formality identifying the Kapranov
$L_\infty$ on $T_X[-1]$ with the Hochschild $L_\infty$ on $\HH^\bullet(\cC)$.
\end{itemize}

\section*{Attack-Heal summary table}

\begin{center}
\begin{tabular}{lll}
Attack & Survived core & Heal \\
\hline
1. $Y_3 = \ell_3^{\min}$? & Yes, after $\Om_X^*$-trace &
  Prop~\ref{prop:y3-l3-equality} \\
2. $\alpha_\BCOV \in H^{0,1}$? & Yes, normalised class &
  Prop~\ref{prop:alpha-bcov-scope} \\
3. Where are the chain-level formulas? & Dolbeault + MC recursion &
  Heal~\ref{heal:3} \\
4. Serre-dual paradox & Serre-dual $\to$ BV-dual & Heal~\ref{heal:4} \\
5. $\hbar$-order assignment & Three different orderings & Heal~\ref{heal:5} \\
6. Birational naturality & Fourier-Mukai intertwining & Heal~\ref{heal:6} \\
7. L-inf on HH${}^\bullet$(C) & Kontsevich formality HKR-ext.\ &
  Heal~\ref{heal:7} \\
8. MC equation recursion & $\ell_n$ unique from MC & Prop~\ref{prop:ln-mc} \\
9. Explicit $Y_3, \alpha_\BCOV$ values & Computed at Q, K3E, CF & Sec~9 \\
10. $R$-matrix relation & Kapranov = leading OPE & Heal~\ref{heal:10} \\
\end{tabular}
\end{center}

\section*{References (working-note abbreviations)}

\begin{itemize}
\item Kapranov 1999, ``Rozansky-Witten invariants via Atiyah classes'',
Compositio Math.\ 115, 71-113.
\item BCOV 1994, ``Kodaira-Spencer theory of gravity and exact results
for quantum string amplitudes'', Comm.\ Math.\ Phys.\ 165, 311.
\item Costello-Li 2016, ``Quantum BCOV theory on Calabi-Yau manifolds
and the higher genus B-model'', arXiv:1201.4501.
\item Kontsevich 1997/2003, ``Deformation quantisation of Poisson
manifolds'', Lett.\ Math.\ Phys.\ 66, 157-216.
\item Markarian 2009, ``The Atiyah class, Hochschild cohomology and the
Riemann-Roch theorem'', J.\ LMS 79, 129-143.
\item C\u ald\u araru 2005, ``The Mukai pairing II'', Adv.\ Math.\ 194,
34-66.
\item Yu 2015, ``A formula for the Atiyah class of a complex of vector
bundles'', Adv.\ Math.\ 271, 303-360.
\item Costello-Francis-Gwilliam 2026, ``The $E_n$-factorisation algebra
of BV observables'' (CFG 2026).
\item Calaque-Rossi 2011, ``Lectures on Duflo isomorphisms in Lie
algebra and complex geometry'', EMS.
\item Li-Ruan 2001; Li 2013: flop-invariance of the tree-level Yukawa.
\item Bridgeland 2002, ``Flops and derived categories'', Invent.\ Math.\
147, 613-632.
\item Candelas-de la Ossa-Green-Parkes 1991, ``A pair of Calabi-Yau
manifolds as an exactly soluble superconformal theory'', Nucl.\ Phys.\
B 359, 21-74.
\item Atiyah 1957, ``Complex analytic connections in fibre bundles'',
Trans.\ AMS 85, 181-207.
\end{itemize}

\end{document}
